\documentclass[11pt,a4paper]{article}
\usepackage[T1]{fontenc}
\usepackage{amsmath}
\usepackage{amssymb}
\usepackage[table]{xcolor}
\usepackage{tabularx}
\usepackage{multirow}
\usepackage{geometry}
\usepackage{enumitem}
\usepackage{parskip}
\usepackage{titlesec}
\usepackage{tocloft}
\usepackage[most]{tcolorbox}

\geometry{margin=2.5cm}

\titleformat{\section}{\large\bfseries}{}{0em}{}
\setlength{\cftbeforesecskip}{6pt}

% Pastel colours, cycled across the questions (q1..q9)
\definecolor{q1color}{HTML}{D6EAF8} % pastel blue
\definecolor{q2color}{HTML}{D5F5E3} % pastel green
\definecolor{q3color}{HTML}{FDEBD0} % pastel orange
\definecolor{q4color}{HTML}{F9E79F} % pastel yellow
\definecolor{q5color}{HTML}{E8DAEF} % pastel purple
\definecolor{q6color}{HTML}{FADBD8} % pastel red/pink
\definecolor{q7color}{HTML}{D1F2EB} % pastel teal
\definecolor{q8color}{HTML}{FCF3CF} % pastel cream
\definecolor{q9color}{HTML}{EBDEF0} % pastel lavender

% Mark-scheme table (FOUR columns: Question | Answer | Marks | Partial Marks)
\newenvironment{mstab}[1]{%
  {\color{#1!75!black}\nobreak Mark Scheme}\par\nopagebreak\smallskip
  \arrayrulecolor{#1!70!black}%
  \renewcommand{\arraystretch}{1.3}%
  \tabularx{\textwidth}{%
    |>{\columncolor{#1!12}\raggedright\arraybackslash}p{1.7cm}|%
     >{\columncolor{#1!12}\raggedright\arraybackslash}p{3.6cm}|%
     >{\columncolor{#1!12}\centering\arraybackslash}p{1.0cm}|%
     >{\columncolor{#1!12}}X|}%
  \hline
  \rowcolor{#1!45} Question & Answer & Marks & Partial Marks \\ \hline
}{%
  \endtabularx
  \arrayrulecolor{black}%
  \par\medskip
}

% Exemplar-answer box
\newtcolorbox{ansbox}[1]{
  colback=#1!8,
  colbacktitle=#1!30,
  colframe=#1!70!black,
  coltitle=black,
  fonttitle=\normalfont,
  title=Worked Solution,
  boxrule=0.6pt,
  arc=2mm,
  left=4pt, right=4pt, top=4pt, bottom=4pt,
  breakable
}

\newcommand{\topiclabel}[1]{\hfill\normalfont\small\itshape\textcolor{black!55}{#1}}

% Per-question head: \examq{paper-id}{number}{topics}{marks}
% Same command and rendered line as the question paper.
\newcommand{\examq}[4]{\section[#1 - Q#2]{#1 - Q#2 - #3 - #4}}

\begin{document}

\begin{center}
  {\LARGE \textbf{Cambridge O Level Mathematics D 4024/12}}\\[6pt]
  {\large May/June 2025 - Paper 1 (Non-calculator): Mark Scheme \& Worked Solutions}
\end{center}

\bigskip

%%%%%%%%%%%%%%%%%%%
\examq{4024/12/M/J/25}{1}{Number Fundamentals \textperiodcentered\ Fractions Fundamentals}{4}
%%%%%%%%%%%%%%%%%%%
\begin{mstab}{q1color}
1(a) & 14 & 1 & \\ \hline
1(b) & 16 & 1 & \\ \hline
1(c) & $\dfrac{12}{45}$ or $\dfrac{4}{15}$ oe & 2 & \textbf{M1} for $\dfrac{2}{9}\times\dfrac{6}{5}$ seen\newline or for $\dfrac{4}{18} \div \dfrac{15}{18}$ seen \\ \hline
\end{mstab}

\begin{ansbox}{q1color}
(a) $6 - 2 \times (-4)$\\[4pt]
Multiplication comes before subtraction, and a positive times a negative is negative:
\[
2 \times (-4) = -8
\]
\[
6 - (-8) = 6 + 8 = \boxed{14}
\]
(b) $80 \times 0.2$\\[4pt]
Note $0.2 = \dfrac{2}{10}$, so
\[
80 \times 0.2 = \frac{80 \times 2}{10} = \frac{160}{10} = \boxed{16}
\]
(c) $\dfrac{2}{9} \div \dfrac{5}{6}$\\[4pt]
To divide by a fraction, multiply by its reciprocal:
\[
\frac{2}{9} \div \frac{5}{6} = \frac{2}{9} \times \frac{6}{5} = \frac{12}{45}
\]
Cancel the common factor 3:
\[
\frac{12}{45} = \boxed{\frac{4}{15}}
\]
\end{ansbox}

%%%%%%%%%%%%%%%%%%%
\examq{4024/12/M/J/25}{2}{Probability Fundamentals}{2}
%%%%%%%%%%%%%%%%%%%
\begin{mstab}{q2color}
2(a) & $\dfrac{4}{11}$ & 1 & \\ \hline
2(b) & $\dfrac{6}{11}$ & 1 & \\ \hline
\end{mstab}

\begin{ansbox}{q2color}
The bag has 11 balls: 5 blue, 4 yellow, and the rest green, so
\[
\text{green} = 11 - 5 - 4 = 2
\]
(a) $P(\text{yellow})$\\[4pt]
There are 4 yellow balls out of 11:
\[
P(\text{yellow}) = \boxed{\frac{4}{11}}
\]
(b) $P(\text{not blue})$\\[4pt]
The balls that are not blue are the 4 yellow and the 2 green, that is $4 + 2 = 6$ balls:
\[
P(\text{not blue}) = \boxed{\frac{6}{11}}
\]
(Equivalently $1 - \dfrac{5}{11} = \dfrac{6}{11}$.)
\end{ansbox}

%%%%%%%%%%%%%%%%%%%
\examq{4024/12/M/J/25}{3}{Geometric Fundamentals}{3}
%%%%%%%%%%%%%%%%%%%
\begin{mstab}{q3color}
3 & Correct net of cube with edge length 3 cm & 3 & \textbf{B2} for net of open cube with edge length 3 cm or for correct net with no/incorrect internal lines\newline OR\newline \textbf{M1} for square with edge length 3 cm soi\newline \textbf{B1} for net of cube seen, any edge length \\ \hline
\end{mstab}

\begin{ansbox}{q3color}
Net of the cube\\[4pt]
Each face of a cube is a square. If the area of one face is $9\,\text{cm}^2$ then the edge length is
\[
\text{edge} = \sqrt{9} = 3\ \text{cm}
\]
On the 1 cm grid, draw an accurate net made of \textbf{six} squares, each $3 \times 3$ grid squares, joined edge to edge so that they fold into a cube. The standard arrangement is a cross: a vertical strip of four squares, with one extra square attached to the left of the second square and one attached to the right of the second square.\\[4pt]
Key requirements: exactly six squares, every square with side $3$ cm, all joined along complete edges, and the internal fold lines drawn correctly (a $2 \times 3$ block of six squares is \textbf{not} a valid net, because it does not fold into a cube).
\end{ansbox}

%%%%%%%%%%%%%%%%%%%
\examq{4024/12/M/J/25}{4}{Angles in Polygons \& Parallel Lines}{3}
%%%%%%%%%%%%%%%%%%%
\begin{mstab}{q4color}
4 & 22 & 3 & \textbf{B2} for $AEB = 110^\circ$ or $BEC = 70^\circ$\newline or \textbf{M2} for $[x =]\ 180 - ((180 - 2\times35) + 48)$ oe\newline or \textbf{B1} for $DCE = 35^\circ$ or $DEC = 110^\circ$ \\ \hline
\end{mstab}

\begin{ansbox}{q4color}
Find the value of $x$\\[4pt]
Triangle $DEC$ has $ED = EC$, so it is isosceles and the base angles at $D$ and $C$ are equal:
\[
\angle DCE = \angle EDC = 35^\circ
\]
Using the angle sum of triangle $DEC$:
\[
\angle DEC = 180^\circ - 2\times35^\circ = 180^\circ - 70^\circ = 110^\circ
\]
$AEC$ and $BED$ are straight lines crossing at $E$, so $\angle AEB$ and $\angle DEC$ are vertically opposite:
\[
\angle AEB = 110^\circ
\]
Now use the angle sum of triangle $AEB$, where $\angle ABE = 48^\circ$:
\[
x = 180^\circ - 110^\circ - 48^\circ = \boxed{22}
\]
\end{ansbox}

%%%%%%%%%%%%%%%%%%%
\examq{4024/12/M/J/25}{5}{Linear Equations \& Inequalities}{2}
%%%%%%%%%%%%%%%%%%%
\begin{mstab}{q5color}
5 & $-3$ & 2 & \textbf{M1} for $4 - x = \dfrac{35}{5}$ or $20 - 5x = 35$ or better \\ \hline
\end{mstab}

\begin{ansbox}{q5color}
Solve $5(4 - x) = 35$\\[4pt]
Divide both sides by 5:
\[
4 - x = \frac{35}{5} = 7
\]
Rearrange:
\[
-x = 7 - 4 = 3 \quad\Longrightarrow\quad x = \boxed{-3}
\]
Check: $5(4 - (-3)) = 5 \times 7 = 35$.\\[4pt]
(Alternatively, expand first: $20 - 5x = 35$, so $-5x = 15$ and $x = -3$.)
\end{ansbox}

%%%%%%%%%%%%%%%%%%%
\examq{4024/12/M/J/25}{6}{Bearings, Constructions \& Scale Drawings}{4}
%%%%%%%%%%%%%%%%%%%
\begin{mstab}{q6color}
6(a) & 29 to 31 & 2 & \textbf{B1} for 5.8 to 6.2 seen\newline or \textbf{M1} for \textit{their} distance in cm $\times$ 5 evaluated \\ \hline
6(b) & $C$ marked in correct position & 2 & \textbf{B1} for bearing of $060^\circ$ from $A$\newline or bearing of $320^\circ$ from $B$ \\ \hline
\end{mstab}

\begin{ansbox}{q6color}
(a) Actual distance from $A$ to $B$\\[4pt]
Measure $AB$ on the scale drawing with a ruler: it is $6\,\text{cm}$ (accept $5.8$ to $6.2$). The scale is 1 cm to 5 km, so
\[
\text{actual distance} = 6 \times 5 = \boxed{30}\ \text{km}
\]
(any value from 29 to 31 is accepted).\\[6pt]
(b) Position of village $C$\\[4pt]
Use a protractor at each village, measuring clockwise from the North line drawn at that village.\\[2pt]
From $A$: measure $060^\circ$ clockwise from $A$'s North line and draw a long ruled line from $A$ in that direction.\\[2pt]
From $B$: measure $320^\circ$ clockwise from $B$'s North line (that is, $40^\circ$ to the \textbf{west} of North) and draw a long ruled line from $B$ in that direction.\\[4pt]
Village $C$ is where the two lines intersect. Mark this point clearly and label it $C$. It lies up and to the right of $A$, and up and to the left of $B$.
\end{ansbox}

%%%%%%%%%%%%%%%%%%%
\examq{4024/12/M/J/25}{7}{Powers, Roots \& Standard Form}{3}
%%%%%%%%%%%%%%%%%%%
\begin{mstab}{q7color}
7(a) & 5 & 1 & \\ \hline
7(b) & $\dfrac{1}{16}$ or 0.0625 & 2 & \textbf{B1} for $\dfrac{1}{4^2}$ seen or $16^{-1}$ seen \\ \hline
\end{mstab}

\begin{ansbox}{q7color}
(a) Evaluate $\sqrt[3]{125}$\\[4pt]
The cube root of 125 is the number whose cube is 125. Since $5^3 = 5\times5\times5 = 125$:
\[
\sqrt[3]{125} = \boxed{5}
\]
(b) Evaluate $4^{-2}$\\[4pt]
A negative index means take the reciprocal of the positive power:
\[
4^{-2} = \frac{1}{4^2} = \frac{1}{16} = \boxed{0.0625}
\]
\end{ansbox}

%%%%%%%%%%%%%%%%%%%
\examq{4024/12/M/J/25}{8}{Scatter Graphs \& Correlation}{4}
%%%%%%%%%%%%%%%%%%%
\begin{mstab}{q8color}
8(a) & 4 points correctly plotted & 2 & \textbf{B1} for 2 correct plots \\ \hline
8(b) & Ruled line of best fit & 1 & \\ \hline
8(c) & Reading at 6.8 km from \textit{their} line of best fit & 1 & Dependent on straight line of best fit with positive gradient \\ \hline
\end{mstab}

\begin{ansbox}{q8color}
(a) Complete the scatter diagram\\[4pt]
Six points are already plotted, so the four remaining pairs from the table must be added:
\[
(4.2,\ 52), \qquad (9.0,\ 116), \qquad (3.8,\ 49), \qquad (5.6,\ 62)
\]
Plot each with a small cross, reading the distance on the horizontal axis and the time on the vertical axis.\\[6pt]
(b) Line of best fit\\[4pt]
The points show strong positive correlation (longer distance means longer time). Draw a single \textbf{ruled} straight line with a \textbf{positive} gradient that passes through the middle of the data, with roughly equal numbers of points above and below it. It should pass close to $(4,\ 52)$ and $(9,\ 114)$.\\[6pt]
(c) Estimate for a 6.8 km walk\\[4pt]
Draw a vertical line up from $6.8$ on the distance axis to the line of best fit, then read horizontally across to the time axis.
\[
\text{Time} \approx \boxed{88}\ \text{minutes}
\]
Any value read correctly from your own line of best fit is accepted (roughly 85 to 90 minutes for a well-drawn line).
\end{ansbox}

%%%%%%%%%%%%%%%%%%%
\examq{4024/12/M/J/25}{9}{Rounding, Estimation \& Bounds}{2}
%%%%%%%%%%%%%%%%%%%
\begin{mstab}{q9color}
9 & 90 and 20 seen as rounded values\newline \textbf{and}\newline final answer 1800 & 2 & \textbf{B1} for 90 and 20 seen as rounded values\newline or \textbf{M1} for product $87.1 \times 23.6$ soi \\ \hline
\end{mstab}

\begin{ansbox}{q9color}
Estimate for the area of the rectangle\\[4pt]
Round each measurement to 1 significant figure:
\[
87.1 \approx 90, \qquad 23.6 \approx 20
\]
Now multiply:
\[
\text{Area} \approx 90 \times 20 = 9 \times 2 \times 100 = \boxed{1800}\ \text{mm}^2
\]
\end{ansbox}

%%%%%%%%%%%%%%%%%%%
\examq{4024/12/M/J/25}{10}{Prime Factors, HCF \& LCM}{3}
%%%%%%%%%%%%%%%%%%%
\begin{mstab}{q1color}
10(a) & $2^2 \times 3 \times 19$\newline or $2 \times 2 \times 3 \times 19$ & 2 & \textbf{B1} for 2, 2, 3, 19\newline or \textbf{M1} for any two stages correct in factor tree or ladder method \\ \hline
10(b) & $2^4 \times 3^2 \times 19^2$\newline or $2\times2\times2\times2\times3\times3\times19\times19$ & 1 & \textbf{FT} \textit{their} product with each term squared \\ \hline
\end{mstab}

\begin{ansbox}{q1color}
(a) 228 as a product of prime factors\\[4pt]
Divide repeatedly by primes:
\[
228 \div 2 = 114, \qquad 114 \div 2 = 57, \qquad 57 \div 3 = 19, \qquad 19 \text{ is prime}
\]
So
\[
228 = 2 \times 2 \times 3 \times 19 = \boxed{2^2 \times 3 \times 19}
\]
(b) 51\,984 as a product of prime factors\\[4pt]
Since $51\,984 = 228^2$, square the factorisation from part (a) by doubling each index:
\[
51\,984 = \left(2^2 \times 3 \times 19\right)^2 = 2^{4} \times 3^{2} \times 19^{2}
\]
\[
\boxed{2^4 \times 3^2 \times 19^2}
\]
\end{ansbox}

%%%%%%%%%%%%%%%%%%%
\examq{4024/12/M/J/25}{11}{Forming \& Solving Equations \textperiodcentered\ Simultaneous Equations}{5}
%%%%%%%%%%%%%%%%%%%
\begin{mstab}{q2color}
11(a) & $4x + 6y = 30$\newline leading to $2x + 3y = 15$ with no errors & 1 & \\ \hline
11(b) & $6x + y = 13$ oe & 1 & \\ \hline
\multirow{2}{*}{11(c)} & Correct method to eliminate one variable & \textbf{M1} & \\ \cline{2-4}
 & $[\text{Small box} =]$ 1.5 oe\newline $[\text{Large box} =]$ 4 & \textbf{A2} & \textbf{A1} for $x = 1.5$ oe or $y = 4$\newline If \textbf{A0} scored, \textbf{SC1} for a pair of values that satisfy either equation \\ \hline
\end{mstab}

\begin{ansbox}{q2color}
(a) Show that $2x + 3y = 15$\\[4pt]
The mass of 4 small boxes is $4x$ and of 6 large boxes is $6y$, and the total is $30$ kg:
\[
4x + 6y = 30
\]
Divide every term by 2:
\[
2x + 3y = 15 \qquad \text{as required.}
\]
(b) Second equation\\[4pt]
6 small boxes and 1 large box have total mass 13 kg:
\[
\boxed{6x + y = 13}
\]
(c) Solve the simultaneous equations\\[4pt]
Make $y$ the subject of the second equation:
\[
y = 13 - 6x
\]
Substitute into $2x + 3y = 15$ to eliminate $y$:
\[
2x + 3(13 - 6x) = 15
\]
\[
2x + 39 - 18x = 15
\]
\[
-16x = 15 - 39 = -24
\]
\[
x = \frac{-24}{-16} = \frac{3}{2} = 1.5
\]
Substitute back:
\[
y = 13 - 6(1.5) = 13 - 9 = 4
\]
\[
\text{Small box} = \boxed{1.5}\ \text{kg}, \qquad \text{Large box} = \boxed{4}\ \text{kg}
\]
Check in the first equation: $2(1.5) + 3(4) = 3 + 12 = 15$.
\end{ansbox}

%%%%%%%%%%%%%%%%%%%
\examq{4024/12/M/J/25}{12}{Transformations}{8}
%%%%%%%%%%%%%%%%%%%
\begin{mstab}{q3color}
12(a) & Rotation\newline $90^\circ$ clockwise oe\newline $(-2,\ 1)$ & 3 & \textbf{B1} for each \\ \hline
12(b) & Triangle at $(-6,\ 1),\ (-3,\ 1),\ (-5,\ 2)$ & 2 & \textbf{B1} for reflection in $x = k$ or in $y = -1$ \\ \hline
12(c)(i) & 3 cao & 1 & \\ \hline
12(c)(ii) & $(3,\ 3)$\newline $(9,\ 6)$ & 2 & \textbf{B1} for $(3,\ 3)$ or $(9,\ 6)$ \\ \hline
\end{mstab}

\begin{ansbox}{q3color}
Triangle $A$ has vertices $(1,\ 1)$, $(4,\ 1)$, $(3,\ 2)$ and triangle $B$ has vertices $(-2,\ -2)$, $(-2,\ -5)$, $(-1,\ -4)$.\\[6pt]
(a) The single transformation\\[4pt]
The triangles are the same size but turned, so the transformation is a rotation. Using tracing paper, or by testing the rule for a $90^\circ$ clockwise turn about $(-2,\ 1)$, namely $(x,\ y) \mapsto \big(-2 + (y - 1),\ 1 - (x + 2)\big)$:
\[
(1,1) \mapsto (-2,\ -2), \quad (4,1) \mapsto (-2,\ -5), \quad (3,2) \mapsto (-1,\ -4)
\]
which is exactly triangle $B$. So the full description is a
\[
\boxed{\text{Rotation, } 90^\circ \text{ clockwise, about } (-2,\ 1)}
\]
All three pieces of information (rotation, angle and direction, centre) are needed.\\[6pt]
(b) Reflection of $A$ in $x = -1$\\[4pt]
Reflecting in the vertical line $x = -1$ keeps the $y$-coordinate and sends $x$ to $-2 - x$:
\[
(1,1) \mapsto (-3,\ 1), \quad (4,1) \mapsto (-6,\ 1), \quad (3,2) \mapsto (-5,\ 2)
\]
Draw the triangle with vertices
\[
\boxed{(-6,\ 1),\ (-3,\ 1),\ (-5,\ 2)}
\]
(c)(i) The scale factor $k$\\[4pt]
The enlargement has centre $(0,\ 0)$, so each vertex is multiplied by $k$. The vertex $(12,\ 3)$ must be the image of $(4,\ 1)$:
\[
k \times 4 = 12 \quad\Longrightarrow\quad k = \boxed{3}
\]
(and $3 \times 1 = 3$, which confirms it).\\[6pt]
(c)(ii) The other two vertices of $C$\\[4pt]
Multiply the other two vertices of $A$ by 3:
\[
3 \times (1,\ 1) = \boxed{(3,\ 3)}, \qquad 3 \times (3,\ 2) = \boxed{(9,\ 6)}
\]
\end{ansbox}

%%%%%%%%%%%%%%%%%%%
\examq{4024/12/M/J/25}{13}{Percentages Fundamentals}{4}
%%%%%%%%%%%%%%%%%%%
\begin{mstab}{q4color}
13(a) & 68 & 2 & \textbf{M1} for $85 - \dfrac{20}{100}\times85$ oe\newline or \textbf{B1} for answer 17 \\ \hline
13(b) & 50 & 2 & \textbf{M1} for $\dfrac{(100 - 20)}{100}x = 40$ oe soi \\ \hline
\end{mstab}

\begin{ansbox}{q4color}
(a) Sale price of the coat\\[4pt]
A reduction of $20\%$ leaves $100\% - 20\% = 80\%$ of the original price. Since $80\% = \dfrac{4}{5}$:
\[
\text{sale price} = \frac{4}{5} \times 85 = 4 \times 17 = \boxed{68}
\]
So the coat costs $\$68$. (The reduction itself is $\frac{1}{5} \times 85 = \$17$, and $85 - 17 = 68$.)\\[6pt]
(b) Original price of the shirt\\[4pt]
The sale price is $80\%$ of the original price $x$:
\[
\frac{80}{100}x = 40
\]
\[
0.8x = 40 \quad\Longrightarrow\quad x = \frac{40}{0.8} = \frac{400}{8} = \boxed{50}
\]
So the shirt cost $\$50$ before the sale. (Check: $20\%$ of $50$ is $10$, and $50 - 10 = 40$.)
\end{ansbox}

%%%%%%%%%%%%%%%%%%%
\examq{4024/12/M/J/25}{14}{Vectors}{9}
%%%%%%%%%%%%%%%%%%%
\begin{mstab}{q5color}
14(a) & $\left[\overrightarrow{AB} = \right] \begin{pmatrix} 10 \\ -4 \end{pmatrix}$ & 2 & \textbf{B1} for $\begin{pmatrix} 10 \\ k \end{pmatrix}$ or $\begin{pmatrix} k \\ -4 \end{pmatrix}$ \\ \hline
14(b) & $(3,\ -3)$ & 2 & \textbf{B1} for each or for $\left[\overrightarrow{OC} = \right] \begin{pmatrix} 3 \\ -3 \end{pmatrix}$ \\ \hline
14(c)(i) & $(-2,\ -1)$ & 2 & \textbf{B1} for one value correct\newline or for correct FT for $x$- or $y$-coordinate of $D$\newline $x = \textit{their } 3 - 0.5 \times \textit{their } 10$ evaluated\newline $y = \textit{their }(-3) - 0.5 \times \textit{their }(-4)$ evaluated\newline or \textbf{M1} for $\left[\overrightarrow{DC} = \right] \begin{pmatrix} (\textit{their } 10) \div 2 \\ (\textit{their } -4) \div 2 \end{pmatrix}$ \\ \hline
14(c)(ii) & $2\sqrt{10}$ cao & 3 & \textbf{M2} for $\left(-4 - \textit{their}(-2)\right)^2 + \left(5 - \textit{their}(-1)\right)^2$ oe\newline or \textbf{M1} for $\left(-4 - \textit{their}(-2)\right)$ and $\left(5 - \textit{their}(-1)\right)$ oe \\ \hline
\end{mstab}

\begin{ansbox}{q5color}
(a) The column vector $\overrightarrow{AB}$\\[4pt]
With $A = (-4,\ 5)$ and $B = (6,\ 1)$, subtract the position vectors:
\[
\overrightarrow{AB} = \begin{pmatrix} 6 - (-4) \\ 1 - 5 \end{pmatrix} = \boxed{\begin{pmatrix} 10 \\ -4 \end{pmatrix}}
\]
(b) Coordinates of $C$\\[4pt]
Add $\overrightarrow{BC}$ to the position of $B$:
\[
C = \begin{pmatrix} 6 \\ 1 \end{pmatrix} + \begin{pmatrix} -3 \\ -4 \end{pmatrix} = \begin{pmatrix} 3 \\ -3 \end{pmatrix}
\]
\[
C = \boxed{(3,\ -3)}
\]
(c)(i) Coordinates of $D$\\[4pt]
$ABCD$ is a trapezium with $AB$ parallel to $DC$ and $AB = 2DC$, so $\overrightarrow{DC}$ is half of $\overrightarrow{AB}$ (and in the same direction, so that the trapezium does not cross itself):
\[
\overrightarrow{DC} = \frac{1}{2}\begin{pmatrix} 10 \\ -4 \end{pmatrix} = \begin{pmatrix} 5 \\ -2 \end{pmatrix}
\]
Since $\overrightarrow{DC} = C - D$, we get $D = C - \overrightarrow{DC}$:
\[
D = \begin{pmatrix} 3 \\ -3 \end{pmatrix} - \begin{pmatrix} 5 \\ -2 \end{pmatrix} = \begin{pmatrix} -2 \\ -1 \end{pmatrix}
\]
\[
D = \boxed{(-2,\ -1)}
\]
(c)(ii) Length of $AD$\\[4pt]
With $A = (-4,\ 5)$ and $D = (-2,\ -1)$:
\[
\overrightarrow{AD} = \begin{pmatrix} -2 - (-4) \\ -1 - 5 \end{pmatrix} = \begin{pmatrix} 2 \\ -6 \end{pmatrix}
\]
By Pythagoras:
\[
AD = \sqrt{2^2 + (-6)^2} = \sqrt{4 + 36} = \sqrt{40}
\]
Simplify the surd using $40 = 4 \times 10$:
\[
AD = \sqrt{4}\times\sqrt{10} = \boxed{2\sqrt{10}}
\]
\end{ansbox}

%%%%%%%%%%%%%%%%%%%
\examq{4024/12/M/J/25}{15}{Surds}{3}
%%%%%%%%%%%%%%%%%%%
\begin{mstab}{q6color}
15(a) & $3\sqrt{7}$ cao & 2 & \textbf{B1} for $5\sqrt{7}$ or $2\sqrt{7}$ \\ \hline
15(b) & $\dfrac{\sqrt{5}}{5}$ cao & 1 & \\ \hline
\end{mstab}

\begin{ansbox}{q6color}
(a) Simplify $\sqrt{175} - \sqrt{28}$\\[4pt]
Take out the largest square factor from each surd:
\[
\sqrt{175} = \sqrt{25 \times 7} = \sqrt{25}\times\sqrt{7} = 5\sqrt{7}
\]
\[
\sqrt{28} = \sqrt{4 \times 7} = \sqrt{4}\times\sqrt{7} = 2\sqrt{7}
\]
Both are multiples of $\sqrt{7}$, so subtract like terms:
\[
5\sqrt{7} - 2\sqrt{7} = \boxed{3\sqrt{7}}
\]
(b) Rationalise the denominator of $\dfrac{1}{\sqrt{5}}$\\[4pt]
Multiply numerator and denominator by $\sqrt{5}$:
\[
\frac{1}{\sqrt{5}} \times \frac{\sqrt{5}}{\sqrt{5}} = \frac{\sqrt{5}}{(\sqrt{5})^2} = \boxed{\frac{\sqrt{5}}{5}}
\]
\end{ansbox}

%%%%%%%%%%%%%%%%%%%
\examq{4024/12/M/J/25}{16}{Cumulative Frequency}{6}
%%%%%%%%%%%%%%%%%%%
\begin{mstab}{q7color}
16(a)(i) & 32 & 1 & \\ \hline
16(a)(ii) & 10 & 2 & \textbf{B1} for $[\text{UQ} =]$ 35 or $[\text{LQ} =]$ 25 \\ \hline
16(a)(iii) & 8 & 2 & \textbf{M1} for 72 seen\newline If 0 scored, \textbf{SC1} for answer 9 \\ \hline
16(b) & Yes, IQR is lower [for journey time from home to work] oe & 1 & Strict FT \textit{their} \textbf{(a)(ii)} \\ \hline
\end{mstab}

\begin{ansbox}{q7color}
There are 80 people, so the cumulative frequency axis runs from 0 to 80.\\[6pt]
(a)(i) Median\\[4pt]
The median is the reading at a cumulative frequency of $\dfrac{80}{2} = 40$. Reading across from 40 to the curve and down to the time axis:
\[
\text{median} = \boxed{32}\ \text{minutes}
\]
(a)(ii) Interquartile range\\[4pt]
The lower quartile is the reading at $\dfrac{80}{4} = 20$ and the upper quartile is the reading at $\dfrac{3 \times 80}{4} = 60$:
\[
\text{LQ} = 25 \ \text{minutes}, \qquad \text{UQ} = 35\ \text{minutes}
\]
\[
\text{IQR} = 35 - 25 = \boxed{10}\ \text{minutes}
\]
(a)(iii) Number with a journey time of 40 minutes or more\\[4pt]
Read \textbf{up} from 40 minutes to the curve: the cumulative frequency is 72, so 72 people took \textbf{less than} 40 minutes. The rest took 40 minutes or more:
\[
80 - 72 = \boxed{8}\ \text{people}
\]
(b) Is Jay correct?\\[4pt]
Consistency is measured by the spread, and the interquartile range is the measure of spread here.
\[
\text{IQR (home to work)} = 10\ \text{minutes} \qquad \text{IQR (work to home)} = 15\ \text{minutes}
\]
\[
\boxed{\text{Yes}} \text{, because the IQR for home to work is lower, so those times are more consistent.}
\]
\end{ansbox}

%%%%%%%%%%%%%%%%%%%
\examq{4024/12/M/J/25}{17}{Rearranging Formulas \textperiodcentered\ Linear Graphs}{5}
%%%%%%%%%%%%%%%%%%%
\begin{mstab}{q8color}
17(a) & $2 - \dfrac{3}{5}x$ oe final answer & 2 & \textbf{M1} for $5y = 10 - 3x$ or for $y + \dfrac{3}{5}x = \dfrac{10}{5}$ or better \\ \hline
17(b) & $y = \dfrac{5}{3}x - 3$ oe & 3 & \textbf{B1} for $[\text{gradient} =]\ \dfrac{5}{3}$ oe or $-\dfrac{1}{\textit{their}\left(-\frac{3}{5}\right)}$ oe\newline \textbf{M1} for substituting $(6,\ 7)$ into $y = (\textit{their } m)x + c$ oe\newline or for $(y - 7) = (\textit{their } m)(x - 6)$ \\ \hline
\end{mstab}

\begin{ansbox}{q8color}
(a) Make $y$ the subject of $5y + 3x = 10$\\[4pt]
Subtract $3x$ from both sides:
\[
5y = 10 - 3x
\]
Divide every term by 5:
\[
y = \frac{10}{5} - \frac{3x}{5} = \boxed{2 - \frac{3}{5}x}
\]
(b) Equation of line $P$\\[4pt]
From part (a) the gradient of $L$ is $-\dfrac{3}{5}$. For perpendicular lines the gradients multiply to $-1$, so the gradient of $P$ is the negative reciprocal:
\[
m_P = -\frac{1}{-\frac{3}{5}} = \frac{5}{3}
\]
Line $P$ passes through $(6,\ 7)$, so substitute into $y = mx + c$:
\[
7 = \frac{5}{3}(6) + c = 10 + c
\]
\[
c = 7 - 10 = -3
\]
\[
\boxed{y = \frac{5}{3}x - 3}
\]
\end{ansbox}

%%%%%%%%%%%%%%%%%%%
\examq{4024/12/M/J/25}{18}{Real-Life Graphs}{5}
%%%%%%%%%%%%%%%%%%%
\begin{mstab}{q9color}
18(a) & Constant speed oe & 1 & \\ \hline
18(b) & $0.25 \times 40 = 10$ & 1 & \\ \hline
18(c) & 220 & 3 & \textbf{M2} for $\dfrac{1}{2}\times10\times(T + 60) = \text{figs } 14$ oe\newline or \textbf{M1} for correct method to find a relevant area under the graph seen \\ \hline
\end{mstab}

\begin{ansbox}{q9color}
(a) Motion between $t = 40$ and $t = 100$\\[4pt]
The graph is horizontal over this interval, so the speed does not change:
\[
\boxed{\text{The cyclist travels at a constant speed (of } 10\ \text{m/s), i.e. zero acceleration.}}
\]
(b) Show that $v = 10$\\[4pt]
On a speed-time graph the acceleration is the gradient. Between $t = 0$ and $t = 40$ the speed rises from 0 to $v$, so
\[
\text{acceleration} = \frac{v - 0}{40 - 0} = 0.25
\]
\[
v = 0.25 \times 40 = 10 \qquad \text{as required.}
\]
(c) Find $T$\\[4pt]
The distance travelled is the area under the graph, which is a trapezium with parallel sides (the two horizontal lengths) equal to $T$ and $100 - 40 = 60$, and height $v = 10$. The total distance is
\[
1.4\ \text{km} = 1400\ \text{m}
\]
\[
\frac{1}{2}\times 10 \times (T + 60) = 1400
\]
\[
5(T + 60) = 1400
\]
\[
T + 60 = 280 \quad\Longrightarrow\quad T = \boxed{220}
\]
\end{ansbox}

%%%%%%%%%%%%%%%%%%%
\examq{4024/12/M/J/25}{19}{Expanding \& Factorising Brackets \textperiodcentered\ Algebraic Fractions}{4}
%%%%%%%%%%%%%%%%%%%
\begin{mstab}{q1color}
19 & $\dfrac{3(x - 2)}{2x + 7}$ or $\dfrac{3x - 6}{2x + 7}$ final answer & 4 & \textbf{B1} for $3(x + 2)(x - 2)$ seen\newline or $(3x - 6)(x + 2)$ seen or $(3x + 6)(x - 2)$ seen\newline \textbf{B2} for $(2x + 7)(x + 2)$ seen\newline or \textbf{B1} for $2x(x + 2) + 7(x + 2)$ seen\newline or for $x(2x + 7) + 2(2x + 7)$ seen\newline or for $(2x + a)(x + b)$ seen where $ab = 14$ or $a + 2b = 11$ \\ \hline
\end{mstab}

\begin{ansbox}{q1color}
Simplify $\dfrac{3x^2 - 12}{2x^2 + 11x + 14}$\\[4pt]
Factorise the numerator: take out the common factor 3, then use the difference of two squares:
\[
3x^2 - 12 = 3(x^2 - 4) = 3(x + 2)(x - 2)
\]
Factorise the denominator. We need $(2x + a)(x + b)$ with $ab = 14$ and $a + 2b = 11$; taking $a = 7$ and $b = 2$ works:
\[
2x^2 + 11x + 14 = 2x^2 + 4x + 7x + 14 = 2x(x + 2) + 7(x + 2) = (2x + 7)(x + 2)
\]
Now cancel the common factor $(x + 2)$:
\[
\frac{3(x + 2)(x - 2)}{(2x + 7)(x + 2)} = \boxed{\frac{3(x - 2)}{2x + 7}}
\]
which may also be written as $\dfrac{3x - 6}{2x + 7}$.
\end{ansbox}

%%%%%%%%%%%%%%%%%%%
\examq{4024/12/M/J/25}{20}{Fractions, Decimals \& Percentages}{4}
%%%%%%%%%%%%%%%%%%%
\begin{mstab}{q2color}
20 & $\dfrac{8}{11}$ cao & 4 & \textbf{B3} for $\dfrac{72}{99}$ or equivalent fraction seen\newline OR\newline \textbf{B2} for $\dfrac{17}{99}$ seen\newline or \textbf{M1} for $17.17\ldots - 0.17\ldots$ oe or for $99x = 17$ oe\newline \textbf{M1} for correct use of common denominator e.g. $\dfrac{17}{99} + \dfrac{55}{99}$ oe\newline FT \textit{their} $\dfrac{17}{99}$ provided denominator $\neq 9$\newline OR\newline \textbf{B2} for $0.7272\ldots$\newline or \textbf{B1} for $0.5555\ldots$ \\ \hline
\end{mstab}

\begin{ansbox}{q2color}
Work out $0.\dot{1}\dot{7} + \dfrac{5}{9}$\\[4pt]
First convert the recurring decimal to a fraction. Let
\[
x = 0.171717\ldots
\]
Two digits recur, so multiply by 100:
\[
100x = 17.171717\ldots
\]
Subtract the first line from the second:
\[
100x - x = 17.1717\ldots - 0.1717\ldots
\]
\[
99x = 17 \quad\Longrightarrow\quad x = \frac{17}{99}
\]
Now add $\dfrac{5}{9}$, using the common denominator 99 (note $\dfrac{5}{9} = \dfrac{55}{99}$):
\[
\frac{17}{99} + \frac{55}{99} = \frac{72}{99}
\]
Simplify by dividing top and bottom by 9:
\[
\frac{72}{99} = \boxed{\frac{8}{11}}
\]
\end{ansbox}

%%%%%%%%%%%%%%%%%%%
\examq{4024/12/M/J/25}{21}{Vectors}{5}
%%%%%%%%%%%%%%%%%%%
\begin{mstab}{q3color}
21(a) & $\left[\overrightarrow{AX} = \right] \dfrac{3}{5}(\mathbf{b} - \mathbf{a})$ or $\dfrac{3}{5}\mathbf{b} - \dfrac{3}{5}\mathbf{a}$\newline final answer & 2 & \textbf{B1} for $\overrightarrow{AB} = \mathbf{b} - \mathbf{a}$ oe \\ \hline
21(b) & $\left[\overrightarrow{XC} = \right] \dfrac{3}{5}\mathbf{a} + \dfrac{9}{10}\mathbf{b}$ final answer & 3 & \textbf{B2} for $\overrightarrow{XC} = \dfrac{3}{2}(\mathbf{a} + \textit{their } AX)$ oe\newline or $\overrightarrow{XC} = -(\textit{their } AX) + \dfrac{3}{2}\mathbf{b}$ oe\newline OR\newline \textbf{M1} for a correct route for $XC$ along lines of diagram\newline \textbf{B1} for $\overrightarrow{AC} = \dfrac{3}{2}\mathbf{b}$ oe or $\overrightarrow{OC} = \mathbf{a} + \dfrac{3}{2}\mathbf{b}$ oe\newline or $\overrightarrow{XC} = \dfrac{3}{2}\overrightarrow{OX}$ oe \\ \hline
\end{mstab}

\begin{ansbox}{q3color}
(a) Find $\overrightarrow{AX}$\\[4pt]
First find $\overrightarrow{AB}$ by going from $A$ back to $O$ and then out to $B$:
\[
\overrightarrow{AB} = -\mathbf{a} + \mathbf{b} = \mathbf{b} - \mathbf{a}
\]
$X$ divides $AB$ in the ratio $AX : XB = 3 : 2$, so $AX$ is $\dfrac{3}{5}$ of $AB$:
\[
\overrightarrow{AX} = \frac{3}{5}(\mathbf{b} - \mathbf{a}) = \boxed{\frac{3}{5}\mathbf{b} - \frac{3}{5}\mathbf{a}}
\]
(b) Find $\overrightarrow{XC}$\\[4pt]
First find $\overrightarrow{OX}$:
\[
\overrightarrow{OX} = \mathbf{a} + \overrightarrow{AX} = \mathbf{a} + \frac{3}{5}\mathbf{b} - \frac{3}{5}\mathbf{a} = \frac{2}{5}\mathbf{a} + \frac{3}{5}\mathbf{b}
\]
$AC$ is parallel to $OB$, so $\overrightarrow{AC} = k\mathbf{b}$ for some $k$, and therefore
\[
\overrightarrow{OC} = \mathbf{a} + k\mathbf{b}
\]
But $O$, $X$ and $C$ are collinear, so $\overrightarrow{OC}$ must be a multiple of $\overrightarrow{OX}$: $\overrightarrow{OC} = \lambda\left(\frac{2}{5}\mathbf{a} + \frac{3}{5}\mathbf{b}\right)$. Comparing the $\mathbf{a}$ components:
\[
\frac{2}{5}\lambda = 1 \quad\Longrightarrow\quad \lambda = \frac{5}{2}
\]
Comparing the $\mathbf{b}$ components:
\[
k = \frac{5}{2}\times\frac{3}{5} = \frac{3}{2}, \qquad \text{so} \qquad \overrightarrow{OC} = \mathbf{a} + \frac{3}{2}\mathbf{b}
\]
Finally:
\[
\overrightarrow{XC} = \overrightarrow{OC} - \overrightarrow{OX} = \left(\mathbf{a} + \frac{3}{2}\mathbf{b}\right) - \left(\frac{2}{5}\mathbf{a} + \frac{3}{5}\mathbf{b}\right)
\]
\[
= \left(1 - \frac{2}{5}\right)\mathbf{a} + \left(\frac{3}{2} - \frac{3}{5}\right)\mathbf{b} = \frac{3}{5}\mathbf{a} + \left(\frac{15 - 6}{10}\right)\mathbf{b}
\]
\[
\boxed{\overrightarrow{XC} = \frac{3}{5}\mathbf{a} + \frac{9}{10}\mathbf{b}}
\]
(Note this equals $\dfrac{3}{2}\overrightarrow{OX}$, as expected since $OXC$ is a straight line.)
\end{ansbox}

%%%%%%%%%%%%%%%%%%%
\examq{4024/12/M/J/25}{22}{Quadratic Equations \textperiodcentered\ Quadratic Graphs}{7}
%%%%%%%%%%%%%%%%%%%
\begin{mstab}{q4color}
22(a) & $(x + 2)^2 - 16$ final answer & 2 & \textbf{B1} for $(x + 2)^2$ seen \\ \hline
22(b) & $(-2,\ -16)$ & 1 & or \textbf{FT} $\left(-(\textit{their } 2),\ \textit{their }(-16)\right)$ \\ \hline
22(c) & Correct sketch with minimum in correct quadrant and values $-6$, 2, $-12$ labelled\newline (mark scheme shows a U-shaped parabola cutting the $x$-axis at $-6$ and 2 and the $y$-axis at $-12$) & 4 & \textbf{B1} for U-shape curve\newline \textbf{B1} for value $-12$ labelled as $y$-intercept\newline \textbf{B2} for values $-6$ and 2 labelled as $x$-intercepts and no extras\newline or \textbf{B1} for parabola with value $-6$ or 2 labelled as an $x$-intercept\newline or \textbf{M1} for $(x + 6)(x - 2)$ soi\newline or for $x + 2 = \pm\sqrt{16}$ FT \textit{their} \textbf{(a)}\newline Max 3 marks if sketch incorrect \\ \hline
\end{mstab}

\begin{ansbox}{q4color}
(a) Complete the square\\[4pt]
Halve the coefficient of $x$ (that is $4 \div 2 = 2$) and use $(x + 2)^2 = x^2 + 4x + 4$:
\[
x^2 + 4x - 12 = (x + 2)^2 - 4 - 12
\]
\[
= \boxed{(x + 2)^2 - 16}
\]
(b) Turning point\\[4pt]
The squared bracket is smallest (equal to 0) when $x + 2 = 0$, that is $x = -2$, and then $y = -16$. So the turning point is a minimum at
\[
\boxed{(-2,\ -16)}
\]
(c) Sketch of $y = x^2 + 4x - 12$\\[4pt]
Find the intercepts. For the $x$-intercepts, factorise (or use part (a)):
\[
x^2 + 4x - 12 = (x + 6)(x - 2) = 0 \quad\Longrightarrow\quad x = -6 \ \text{ or } \ x = 2
\]
For the $y$-intercept, put $x = 0$:
\[
y = 0 + 0 - 12 = -12
\]
Draw a \textbf{U-shaped} parabola (positive $x^2$ coefficient) that
\[
\text{crosses the } x\text{-axis at } \boxed{-6} \text{ and } \boxed{2}, \qquad \text{crosses the } y\text{-axis at } \boxed{-12},
\]
with its minimum at $(-2,\ -16)$, which lies below the $x$-axis and to the left of the $y$-axis (the third quadrant). Label the values $-6$, $2$ and $-12$ on the axes, and label no other intercepts.
\end{ansbox}

%%%%%%%%%%%%%%%%%%%
\examq{4024/12/M/J/25}{23}{Circles, Arcs \& Sectors}{5}
%%%%%%%%%%%%%%%%%%%
\begin{mstab}{q5color}
23 & $26\pi$ cao & 5 & \textbf{M2} for $AOB = 100$\newline or $AOB = \dfrac{5\pi \times 360}{2\pi \times 9}$ or better\newline or \textbf{M1} for $\dfrac{x}{360}\times 2\pi \times 9 = 5\pi$ oe\newline \textbf{M2} for $\dfrac{360 - \textit{their } AOB}{360}\times\pi\times6^2$ oe\newline or $\pi\times6^2 - \dfrac{\textit{their } AOB}{360}\times\pi\times6^2$ oe\newline or \textbf{M1} for $\dfrac{\textit{their } AOB}{360}\times\pi\times6^2$\newline or for major sector angle $= 360 - \textit{their } AOB$ \\ \hline
\end{mstab}

\begin{ansbox}{q5color}
Area of the major sector $OCD$\\[4pt]
\textbf{Step 1: find the sector angle $AOB$.}\\
The minor sector $OAB$ has radius $OA = 9$ and arc length $5\pi$. Let the angle $AOB$ be $\theta$:
\[
\frac{\theta}{360}\times 2\pi \times 9 = 5\pi
\]
\[
\frac{\theta}{360}\times 18\pi = 5\pi \quad\Longrightarrow\quad \frac{\theta}{360} = \frac{5}{18}
\]
\[
\theta = \frac{5}{18}\times 360 = 5 \times 20 = 100^\circ
\]
\textbf{Step 2: find the angle of the major sector.}\\
$OCA$ and $ODB$ are straight lines, so the major sector $OCD$ is the rest of the full turn about $O$:
\[
\text{angle } COD = 360^\circ - 100^\circ = 260^\circ
\]
\textbf{Step 3: find its area.}\\
The major sector has radius $OC = 6$:
\[
\text{Area} = \frac{260}{360}\times \pi \times 6^2 = \frac{13}{18}\times 36\pi
\]
\[
= 13 \times 2\pi = \boxed{26\pi}\ \text{cm}^2
\]
\end{ansbox}

\end{document}